\documentclass[12pt]{beamer}
% \usetheme{umbc4}
\usetheme{moloch}

\usepackage{amsmath,amscd}
\usepackage{fontspec}
\usepackage{unicode-math}
\usepackage{pifont}
\usepackage{xcolor}
\usepackage[dvipsnames]{xcolor}
\usepackage{listings}

\definecolor{midgray}{rgb}{0.65, 0.65, 0.65}

% Set the slide body to black background with white text
\setbeamercolor{background canvas}{bg=black}
\setbeamercolor{normal text}{fg=white}

% Set the title bar to black background with white text
\setbeamercolor{frametitle}{bg=black, fg=white}

% Ensure other structural elements (bullets, etc.) are visible
\setbeamercolor{structure}{fg=white}

\setmainfont{Libertinus Serif}
\setsansfont{Libertinus Sans}
\setmathfont{STIX Two Math}

\newcommand{\FrameTitle}[2]{\frametitle{#1 \hfill \small\textnormal{\textcolor{midgray}{#2}}}}
\newcommand{\snl}{\vspace*{3pt}\\}
\newcommand{\nl}{\vspace*{1em}\\}
\newcommand{\vsp}{\vspace*{1em}}
\newcommand{\vspthreept}{\vspace*{3pt}}

\definecolor{VeyLightBlue}{rgb}{0.65, 0.8, 1.0}
\definecolor{LightGreen}{rgb}{0.65, 1.0, 0.8}
\definecolor{IceBlue}{rgb}{0.6, 0.9, 0.93}

\begin{document}
\title{\center{Basic Mathematical Definitions \\ \textit{and} \\ Using Python for Cryptography \\}}
\author{\center{Manipal School of Information Sciences \\ Manipal \\}}
\date{\center{Friday, 24 April 2026}}
\maketitle

% Set up customized Python listing style
\lstdefinestyle{pythonstyle}{
    language=Python,
    backgroundcolor=\color{black},
    commentstyle=\color{LightGreen},
    basicstyle=\ttfamily\footnotesize\color{white},
    keywordstyle=\color{VeyLightBlue},
    stringstyle=\color{orange},
    breaklines=false,
    linewidth=\linewidth,
    xleftmargin=8pt,
    framexleftmargin=8pt,
    numbers=left,           % Position of line numbers
    numberstyle=\tiny,      % Size of line numbers
    stepnumber=1,           % Numbering interval
    numbersep=2pt,          % Distance from code
    gobble=6,              % Remove leading spaces (adjust as needed)
    rulecolor=\color{gray},
    tabsize=2,               % Reduces Tab width
    columns=fullflexible,        % Compresses character spacing
    keepspaces=true,         % Ensures the indentation stays consistent
    frame=lr,
    framerule=0.5pt,          % Adjust thickness
}
\lstset{style=pythonstyle}

\begin{frame}\FrameTitle{Rule of Exponentiation}{Math101}
    {\huge
    $$
        {g^{a^{b}} = g^{b^{a}} = g^{ab}}
    $$
    }
\end{frame}

\begin{frame}[fragile]\FrameTitle{Exponentiation}{Math101}
    \begin{lstlisting}
        def exp():
            g = 9 # generator
            a = 13 # a's secret key (aka private key)
            a_public = pow(g, a) # g^a
            b = 42 # b's secret key
            b_public = pow(g, b) # g^b
            # @a shares their public key `a_public` with @b
            ss_b = pow(a_public, b_secret)
            # @b shares their public key `b_public` with @a
            ss_a = pow(b_public, a_secret)
            assert ss_a == ss_b
    \end{lstlisting}

    {\Large
    $$
        g^{a^{b}} = g^{b^{a}} = g^{ab}
    $$
    }

\end{frame}

\begin{frame}[fragile]\FrameTitle{Modular Exponentiation}{Math101}
    \begin{lstlisting}
        def mod_exp():
            p = 3329 # prime modulus
            g = 9 # generator
            a_secret = 13
            a_public = pow(g, a_secret, p) # g^a_secret mod p
            b_secret = 42
            b_public = pow(g, b_secret, p) # g^b_secret mod p
            ss_alice = pow(b_public, a_secret, p)
            ss_bob = pow(a_public, b_secret, p)
            assert ss_alice == ss_bob
    \end{lstlisting}

    {\Large
    $$
        g^{a^{b}} \equiv g^{b^{a}} \equiv g^{ab}  \pmod{p}
    $$
    }
\end{frame}

\begin{frame}\FrameTitle{Discrete Logarithm}{Math101}
    {\huge
    Given
    \center{$g$, $p$, and \color{LightGreen}{$g^{\color{Apricot}{a}} \pmod{p}$ },} \nl
    find \color{Apricot}{$a$}
    }
\end{frame}

\begin{frame}\FrameTitle{Definitions - Function and its friends}{Math101}
    \begin{itemize}
        \item[] Function
        \item[] Domain
        \item[] Codomain
        \item[] Range
    \end{itemize}

    \textbf{Function} is a mapping from a \textit{domain} to a \textit{codomain}.\nl
    \textbf{Range} is the subset of the codomain that is \textit{actually} mapped to by the function.
\end{frame}

\begin{frame}\FrameTitle{Definitions - Image and Preimage}{Math101}
    Image of a function $f$ is the set of all outputs of $f$.\nl
    Preimage of a function $f$ is the set of all inputs that map to a given output of $f$.
\end{frame}

\begin{frame}\FrameTitle{Definitions - Second Preimage}{Math101}
    Given a function $f$ and an input $x$, a \textbf{second preimage} of $x$ is an input $x' \neq x$ such that $f(x) = f(x')$.
\end{frame}

\begin{frame}\FrameTitle{Definitions - One Way Function}{Math101}
    A \textbf{one-way function} is a function that is easy to compute but hard to invert.\nl
    \vspace*{1em}
    More formally, a function $f$ is a one-way function if:
    \begin{itemize}
        \item $f$ can be computed in polynomial time.
        \item For every probabilistic polynomial-time algorithm $A$, the probability that $A(f(x))$
            outputs a preimage of $f(x)$ is negligible.
    \end{itemize}
\end{frame}

\begin{frame}\FrameTitle{Hash Functions}{Math101}
    A \textbf{hash function} is a function that maps data of arbitrary size to a fixed-size value.\nl
    \vspace*{1em}
    A \textbf{cryptographic hash function} is a hash function that has the following properties:
    \begin{itemize}
        \item It is easy to compute the hash value for any given input.
        \item It is infeasible to generate a valid input that has a given hash value.
        \item It is infeasible to find two different inputs that produce the same hash value.
    \end{itemize}
\end{frame}

\begin{frame}\FrameTitle{Random Number Generator}{Math101}
    A \textbf{cryptographically secure pseudorandom number generator} (CSPRNG)
    is a pseudorandom number generator that has the following properties:
    \begin{itemize}
        \item It is computationally indistinguishable from a true random number generator.
        \item It is resistant to state compromise extensions, meaning that even if an attacker
                learns the internal state of the generator at some point, they cannot predict future outputs.
    \end{itemize}
\end{frame}
\end{document}
